\section{Outline}
\label{sec:outline}
The structure of this thesis is organized as follows:
\begin{itemize}
  \item \textbf{Chapter 2: Background and Theoretical Foundations} -- A review of the technical foundations of LoRaWAN networks, indoor radio propagation and path loss, machine learning concepts for regression, and TinyML constraints.
  \item \textbf{Chapter 3: Related Work} -- A review of existing literature on federated learning for IoT and resource-constrained environments, including a detailed comparison positioning this work against key reference frameworks.
  \item \textbf{Chapter 4: Methodology} -- A detailed description of the custom system design and architecture, the data processing and cleaning pipeline of the Hölderlinstraße campus dataset, the federated learning protocol, and the edge firmware implementation.
  \item \textbf{Chapter 5: Results and Discussion} -- An in-depth evaluation of the simulation experiments comparing FedAvg, FedProx, SCAFFOLD, and FedYogi optimization algorithms under non-IID conditions, alongside communication efficiency and payload analyses.
  \item \textbf{Chapter 6: Conclusion and Future Work} -- A summary of the key research findings, the positioning of the decentralized approach, and recommendations for future physical edge deployments.
\end{itemize}

The goal of this work is to advance the deployment of decentralized edge intelligence in low-power wide-area networks by offering a communication-efficient federated learning framework that protects data privacy while maintaining robust path loss modeling accuracy.
